In the blank square given in the answer sheet (field of view $45^\circ$,
North direction marked), you should assume that Nova Sagittarii 2 is in the
centre (marked by a green star):
\begin{center}
RA (J2000): $18^{\mathrm{h}} 36^{\mathrm{m}} 56.84^{\mathrm{s}}$ \qquad
Dec (J2000): $-28^\circ 55' 39.8''$
\end{center}

\begin{description}
\item[(1)] Draw the declination and right ascension lines passing through
  Nova Sagittarii. Also draw parallel gridlines at a convenient interval and
  label them accordingly. (The field of view is $45^\circ$.)
  \hfill(10 points)
\item[(2)] Mark with a filled circle the five brightest stars in the field
  of view on the answer sheet, along with their names. You can add several
  other stars for your orientation. \hfill(25 points)
\item[(3)] Draw an unfilled circle representing the Moon's position tomorrow
  night (on 29 July 2015). \hfill(10 points)
\item[(4)] Mark the position of any five Messier objects in the field of
  view by a cross ($\times$), and write down their names beside their marks.
  \hfill(25 points)
\item[(5)] Draw a line representing the galactic plane. \hfill(10 points)
\item[(6)] Mark the positions of Nunki ($\sigma$ Sgr) and Shaula
  ($\lambda$ Sco) by a star ($*$). Draw a line connecting the two stars.
  Estimate the angular distance between these stars and write it on the line
  connecting both stars. \hfill(20 points)
\end{description}

% FIGURE NEEDED: blank square answer chart, field of view 45 degrees, North
% direction marked, with Nova Sagittarii 2 (green star) at the centre
